\section{Prospective Validation Extension}
\subsection{Leakage prevention}
The target is the scheduled child report that the preceding task block was too difficult. Predictor timestamps must precede the presentation of that check-in. The current response, subsequent task performance, and subsequent controller actions are prohibited as inputs. Both models use identical eligible prediction occasions.

\subsection{Participant-level evaluation}
A non-camera baseline model will be compared with an otherwise identical model augmented by local facial movement summaries. Outer evaluation splits are grouped by child. All scaling, feature selection, hyperparameter optimization, and calibration occur inside training participants. The primary comparison remains the equal-child-weighted paired log-loss difference defined earlier. Coverage and missing reports will be tabulated against all scheduled check-ins.

\subsection{Decision utility}
A later randomized controller study is conditional on a reproducible improvement in prediction and acceptable tracking coverage. It will compare camera-enabled decisions with the standard controller while preserving the same safety overrides. Excessive challenge exposure, unnecessary reductions in task difficulty, voluntary stops, and recovery outcomes are prospective decision outcomes.

\subsection{Falsification}
A camera-specific benefit is unsupported when gains disappear in participant-held-out evaluation, are explained by tracking availability, or fail to improve controller decisions. Facial behavior is not interpreted as a categorical emotion label or an ADHD diagnostic signal.